%==============================================================================
% Relatorio Tecnico - Simulador de NoC SPIN (Fat-Tree Quaternaria) em SystemC
% Formato: Artigo IEEE, 2 colunas, <= 8 paginas
% Compilacao recomendada: pdflatex (2x) relatorio.tex   ou   Overleaf
%==============================================================================
\documentclass[conference]{IEEEtran}
\IEEEoverridecommandlockouts

% ----------------------------- Pacotes --------------------------------------
\usepackage[utf8]{inputenc}
\usepackage[T1]{fontenc}
\usepackage[brazil]{babel}
\usepackage{graphicx}
\usepackage{amsmath,amssymb}
\usepackage{siunitx}
\usepackage{booktabs}
\usepackage{array}
\usepackage{listings}
\usepackage{xcolor}
\usepackage{tikz}
\usepackage{url}
\usetikzlibrary{positioning,arrows.meta,calc,fit,backgrounds,shapes.geometric}

% --------------------- Estilo para logs / codigo ----------------------------
\definecolor{logbg}{RGB}{245,245,245}
\definecolor{kw}{RGB}{0,90,160}
\definecolor{cmt}{RGB}{0,120,0}
\definecolor{str}{RGB}{160,40,40}

\lstdefinestyle{cpp}{
  language=C++,
  basicstyle=\ttfamily\scriptsize,
  keywordstyle=\color{kw}\bfseries,
  commentstyle=\color{cmt}\itshape,
  stringstyle=\color{str},
  backgroundcolor=\color{logbg},
  breaklines=true,
  showstringspaces=false,
  columns=fullflexible,
  frame=single,
  framesep=3pt,
  numbers=none,
  tabsize=2,
}
\lstdefinestyle{log}{
  basicstyle=\ttfamily\tiny,
  backgroundcolor=\color{logbg},
  breaklines=true,
  columns=fullflexible,
  frame=single,
  framesep=3pt,
  literate={->}{{$\rightarrow$}}1 {>>}{{>>}}2 {<<}{{<<}}2,
}

\begin{document}

% ============================== TITULO ======================================
\title{Simulador \emph{Cycle-Accurate} de uma Rede-em-Chip SPIN com Topologia
Fat-Tree Quaternária em SystemC: Comutação \emph{Wormhole}, Controle de Fluxo
por Créditos e Arbitragem Probabilística}

\author{
\IEEEauthorblockN{Breno Barbosa de Oliveira \quad
Felipe Marley de Oliveira Gomes \\
Fernando Gomes de Almeida \quad
Gustavo de Freitas Rodrigues}
\IEEEauthorblockA{Curso de Ciência da Computação\\
Departamento de Informática e Matemática Aplicada (DIMAp)\\
Universidade Federal do Rio Grande do Norte (UFRN)\\
Trabalho Prático 2 --- Redes-em-Chip (NoC)}
}

\maketitle

% ============================== RESUMO ======================================
\begin{abstract}
Este trabalho apresenta o projeto, a modelagem e a verificação de um simulador
de ciclo exato (\emph{cycle-accurate}) de uma Rede-em-Chip (\emph{Network-on-Chip}
--- NoC) inspirada na arquitetura SPIN (\emph{Scalable Programmable Integrated
Network}), descrito em SystemC ao nível de transferência entre registradores
(RTL). A rede adota uma topologia em Árvore Gorda (\emph{Fat-Tree}) quaternária
de dois níveis, com 4 roteadores-folha, 4 roteadores-raiz e 16 terminais de
processamento. São implementados os três mecanismos fundamentais de uma NoC
moderna: comutação \emph{wormhole}, na qual mensagens são fragmentadas em
\emph{flits}; controle de fluxo baseado em créditos, que garante ausência de
perda de pacotes por estouro de buffer; e um árbitro probabilístico híbrido,
derivado das especificações da rede SPIN, que sorteia a cada ciclo a ordem de
atendimento das portas para mitigar \emph{starvation} e congestionamento. O
roteamento combina uma fase ascendente adaptativa (escolha do enlace com maior
número de créditos livres) com uma fase descendente determinística (endereçamento
hierárquico pelo identificador do destino). Os resultados de simulação, obtidos
pela injeção de uma mensagem \emph{wormhole} de quatro \emph{flits} entre dois
terminais em sub-árvores distintas, demonstram a correta alocação e liberação de
canais, a integridade do controle de fluxo e a remontagem completa do pacote no
destino, validando funcionalmente o modelo proposto.
\end{abstract}

\begin{IEEEkeywords}
Rede-em-Chip, NoC, SPIN, Fat-Tree, SystemC, comutação wormhole, controle de
fluxo por créditos, roteamento adaptativo, arbitragem probabilística.
\end{IEEEkeywords}

% ============================ INTRODUCAO ====================================
\section{Introdução}

A evolução dos sistemas integrados em chip (\emph{Systems-on-Chip} --- SoC) levou
à integração de dezenas a centenas de núcleos de processamento, memórias e
aceleradores em um mesmo substrato de silício. Nesse cenário, os tradicionais
barramentos compartilhados e as conexões ponto-a-ponto dedicadas tornaram-se
gargalos: o barramento não escala em largura de banda nem em consumo, e as
conexões dedicadas explodem em área de roteamento à medida que o número de
núcleos cresce~\cite{benini2002networks}. As Redes-em-Chip (NoCs) surgem como o
paradigma de comunicação que substitui essas estruturas por uma infraestrutura
de roteadores e enlaces comutados por pacotes, à semelhança das redes de
computadores, porém com restrições severas de área, latência e energia.

Entre as topologias propostas, a Árvore Gorda (\emph{Fat-Tree}) destaca-se por
oferecer largura de banda crescente em direção à raiz e múltiplos caminhos
redundantes entre quaisquer dois nós, propriedades que favorecem o balanceamento
de carga e a tolerância a congestionamento. A rede SPIN~\cite{guerrier2000spin},
uma das primeiras NoCs propostas na literatura, é justamente um exemplo canônico
de NoC baseada em \emph{fat-tree}, empregando roteamento adaptativo na subida e
determinístico na descida.

Este trabalho tem por objetivo projetar, implementar e verificar um simulador
\emph{cycle-accurate} dessa arquitetura utilizando a biblioteca SystemC. Mais do
que reproduzir a topologia, busca-se modelar fielmente, ao nível RTL, os
mecanismos de hardware que governam o tráfego: a fragmentação de mensagens em
\emph{flits} (comutação \emph{wormhole}), o controle de fluxo por créditos entre
roteadores vizinhos e a política de arbitragem probabilística inspirada na SPIN.
As contribuições do trabalho são: (i) um modelo modular e sintetizável-em-espírito
da rede, com separação estrita entre interface e implementação; (ii) a
implementação de uma máquina de estados \emph{wormhole} com alocação e liberação
explícita de canais por meio de uma tabela de travas; e (iii) a validação
funcional do conjunto por meio de simulação dirigida por estímulos.

O restante do artigo organiza-se da seguinte forma. A
Seção~\ref{sec:teoria} apresenta o referencial teórico sobre NoCs, topologia
\emph{fat-tree}, comutação \emph{wormhole}, controle de fluxo e a rede SPIN. A
Seção~\ref{sec:modelo} detalha o modelo implementado, seus componentes e o
diagrama de blocos. A Seção~\ref{sec:resultados} discute as simulações realizadas
e seus resultados. A Seção~\ref{sec:conclusao} conclui o trabalho.

% ========================= REFERENCIAL TEORICO ==============================
\section{Referencial Teórico}
\label{sec:teoria}

\subsection{Redes-em-Chip (NoC)}
Uma NoC é composta por três elementos básicos: \emph{roteadores}, responsáveis
por encaminhar pacotes; \emph{enlaces} (\emph{links}), que conectam roteadores
entre si e a núcleos; e \emph{interfaces de rede} (aqui denominadas terminais),
que adaptam os núcleos ao protocolo da rede. O desempenho de uma NoC é
caracterizado por três eixos de projeto fortemente acoplados: a \emph{topologia}
(o grafo de interconexão), o \emph{algoritmo de roteamento} (a escolha do
caminho) e o \emph{mecanismo de controle de fluxo} (a gestão dos recursos de
buffer e enlace)~\cite{dally2004principles}.

\subsection{Topologia Fat-Tree Quaternária}
A \emph{fat-tree} é uma árvore na qual a capacidade dos enlaces aumenta em
direção à raiz, eliminando o gargalo de banda típico de árvores convencionais.
Em uma \emph{fat-tree} quaternária (aridade 4), cada roteador possui quatro
portas descendentes (em direção às folhas) e quatro ascendentes (em direção à
raiz). A propriedade essencial é a existência de \emph{múltiplos caminhos}
ascendentes equivalentes entre uma folha e a raiz: qualquer roteador-raiz pode
servir de ponto de inflexão para um pacote, o que habilita o roteamento
adaptativo e o balanceamento de carga.

\subsection{Comutação Wormhole}
Na comutação \emph{wormhole}, uma mensagem é dividida em unidades de controle de
fluxo chamadas \emph{flits} (\emph{flow control units}). O primeiro,
\emph{flit-cabeçalho} (\emph{header}), carrega o endereço de destino e, ao
percorrer a rede, \emph{aloca} (tranca) os canais de saída que utiliza,
abrindo um ``túnel''. Os \emph{flits} de corpo (\emph{payload}) seguem o
cabeçalho em fila indiana sem necessidade de re-roteamento, e o \emph{flit-cauda}
(\emph{tail}) libera os canais à medida que avança. A vantagem é a baixa latência
e o reduzido requisito de \emph{buffer} (basta espaço para alguns \emph{flits},
não para o pacote inteiro); a desvantagem é que um pacote bloqueado pode reter
canais ao longo de vários roteadores, exigindo cuidado com \emph{deadlock} e
\emph{starvation}~\cite{dally2004principles}.

\subsection{Controle de Fluxo por Créditos}
O controle de fluxo \emph{credit-based} previne o estouro dos \emph{buffers} de
entrada. Cada transmissor mantém um contador de \emph{créditos} que reflete o
número de posições livres no \emph{buffer} FIFO do receptor a jusante. Um
\emph{flit} só é enviado se houver crédito disponível ($\text{créditos} > 0$);
ao enviar, o transmissor decrementa o contador. Quando o receptor consome um
\emph{flit} de seu \emph{buffer}, devolve um \emph{crédito} ao transmissor, que
incrementa o contador. Esse protocolo garante, por construção, a ausência de
perda de \emph{flits} por transbordo.

\subsection{A Rede SPIN e a Arbitragem}
A SPIN~\cite{guerrier2000spin} é uma NoC \emph{fat-tree} que populariza o esquema
de roteamento em duas fases: \emph{subida adaptativa}, em que o pacote pode tomar
qualquer enlace ascendente livre (tipicamente o menos congestionado), e
\emph{descida determinística}, em que, a partir do ponto de inflexão, o caminho
até o destino é único e calculado pelo endereço hierárquico. Quando múltiplas
entradas disputam a mesma saída, um \emph{árbitro} decide a ordem de atendimento.
Políticas puramente fixas (prioridade estática) podem causar \emph{starvation};
por isso, adota-se aqui uma arbitragem \emph{probabilística}, que sorteia a cada
ciclo a ordem de varredura das portas, garantindo que toda entrada seja
eventualmente atendida.

% ========================= MODELO IMPLEMENTADO ==============================
\section{Modelo Implementado}
\label{sec:modelo}

\subsection{Visão Geral e Tecnologia}
O simulador foi inteiramente descrito em \textbf{SystemC} (Accellera 3.0.2,
padrão C++17), uma biblioteca de classes C++ para modelagem de hardware. A
escolha do SystemC permite descrever o paralelismo intrínseco do hardware por
meio de processos (\texttt{SC\_METHOD}) sensíveis à borda de subida do relógio,
canais (\texttt{sc\_signal}) e portas (\texttt{sc\_in}/\texttt{sc\_out}),
obtendo um modelo de \emph{ciclo exato}. O relógio do sistema opera a um período
de \SI{10}{ns} (\SI{100}{MHz}), conforme definido no \emph{testbench}.

O projeto adota separação estrita entre interface (\texttt{include/spin/*.hpp}) e
implementação (\texttt{src/*.cpp}), com o código de simulação organizado como uma
biblioteca reutilizável e os cenários de exercício isolados em \texttt{tests/}.
A Tabela~\ref{tab:arquivos} resume os módulos.

\begin{table}[!t]
\caption{Organização dos módulos do simulador}
\label{tab:arquivos}
\centering
\renewcommand{\arraystretch}{1.2}
\begin{tabular}{@{}p{1.7cm}p{5.8cm}@{}}
\toprule
\textbf{Módulo} & \textbf{Responsabilidade} \\
\midrule
\texttt{flit}      & Define o tipo \texttt{Flit} e o \texttt{enum class FlitType} (\texttt{kHeader}, \texttt{kPayload}, \texttt{kTail}). \\
\texttt{terminal}  & Interface de rede: injeta a sequência \emph{wormhole} e consome/remonta \emph{flits} recebidos. \\
\texttt{router}    & Roteador RSPIN: \emph{buffers}, créditos, árbitro probabilístico e máquina \emph{wormhole}. \\
\texttt{fat\_tree} & Módulo \emph{top-level}: instancia roteadores e terminais e realiza o cabeamento da malha. \\
\texttt{trace}     & Registro de log opcional (ativável por variável de ambiente). \\
\texttt{tests/}    & Conjunto de \emph{testbenches} \texttt{sc\_main}: cada um injeta um cenário e verifica a entrega. \\
\bottomrule
\end{tabular}
\end{table}

\subsection{Estrutura do Flit}
A unidade de transporte é a \texttt{struct Flit}, que carrega o tipo do
\emph{flit}, os identificadores lógicos de origem (\texttt{src}) e destino
(\texttt{dest}) em campos \texttt{sc\_uint<10>}, uma carga útil textual
(\texttt{data}) usada para inspeção da simulação, e um bit \texttt{valid} que
sinaliza presença de dado no enlace naquele ciclo.

\subsection{Topologia e Diagrama de Blocos}
A rede contém 16 terminais ($T_0 \ldots T_{15}$), 4 roteadores-folha
($L_0 \ldots L_3$, nível 1) e 4 roteadores-raiz ($R_0 \ldots R_3$, nível 2).
Cada roteador-folha conecta-se a 4 terminais por suas portas descendentes e aos
4 roteadores-raiz por suas portas ascendentes, formando a malha
\emph{fat-tree} completamente conectada entre os dois níveis
(Fig.~\ref{fig:topologia}). O terminal $i$ pertence ao roteador-folha
$\lfloor i/4 \rfloor$, na porta $i \bmod 4$.

\begin{figure}[!t]
\centering
\resizebox{\columnwidth}{!}{%
\begin{tikzpicture}[
  every node/.style={font=\scriptsize},
  root/.style={draw,fill=blue!12,minimum width=0.85cm,minimum height=0.6cm,rounded corners=2pt},
  leaf/.style={draw,fill=green!12,minimum width=0.85cm,minimum height=0.6cm,rounded corners=2pt},
  term/.style={draw,fill=orange!18,minimum width=0.5cm,minimum height=0.42cm,font=\tiny,inner sep=1pt},
  ]
  \def\ls{2.4}   % espacamento horizontal entre folhas/raizes
  \def\ts{0.58}  % espacamento horizontal entre terminais de uma mesma folha
  % Raizes (nivel 2)
  \foreach \i in {0,1,2,3} {
    \node[root] (R\i) at (\ls*\i,3.2) {$R_\i$};
  }
  % Folhas (nivel 1)
  \foreach \i in {0,1,2,3} {
    \node[leaf] (L\i) at (\ls*\i,1.6) {$L_\i$};
  }
  % Malha folha-raiz (fat-tree completa)
  \foreach \l in {0,1,2,3} {
    \foreach \r in {0,1,2,3} {
      \draw[gray!55,line width=0.3pt] (L\l.north) -- (R\r.south);
    }
  }
  % Terminais
  \foreach \l in {0,1,2,3} {
    \foreach \p in {0,1,2,3} {
      \pgfmathtruncatemacro{\tid}{4*\l+\p}
      \pgfmathsetmacro{\xx}{\ls*\l + \ts*(\p-1.5)}
      \node[term] (T\tid) at (\xx,0.1) {\tid};
      \draw[gray!70,line width=0.3pt] (L\l.south) -- (T\tid.north);
    }
  }
  % Rotulos de nivel: bem a esquerda, ancorados a leste (texto cresce para fora)
  \node[font=\scriptsize,anchor=east] at (-1.5,3.2) {Nível 2 (raiz)};
  \node[font=\scriptsize,anchor=east] at (-1.5,1.6) {Nível 1 (folha)};
  \node[font=\scriptsize,anchor=east] at (-1.5,0.1) {Terminais};
\end{tikzpicture}%
}
\caption{Topologia \emph{fat-tree} quaternária: 4 roteadores-raiz, 4
roteadores-folha e 16 terminais. Todo enlace folha--raiz transporta dados em um
sentido e créditos no sentido oposto.}
\label{fig:topologia}
\end{figure}

\subsection{O Roteador RSPIN}
O coração do modelo é o módulo \texttt{RSPIN\_Router}, parametrizado por
identificador (\texttt{my\_id}) e nível (\texttt{level}). Cada roteador possui
quatro portas ascendentes e quatro descendentes, cada uma com seu par de sinais
de dados e de crédito. Internamente mantém: \emph{buffers} FIFO de entrada
(\texttt{std::queue}) de profundidade 4 por porta; contadores de crédito de
saída; \emph{flags} de ocupação de saída (\texttt{out\_*\_busy}); e a tabela
\texttt{input\_locks[8]}, que registra, para cada uma das 8 entradas, qual
direção e porta de saída ela ``trancou''. O processo \texttt{routing\_logic} é
disparado a cada borda de relógio e executa o pipeline lógico da
Fig.~\ref{fig:fsm}.

\begin{figure}[!t]
\centering
\begin{tikzpicture}[
  node distance=4.5mm,
  every node/.style={font=\scriptsize},
  stage/.style={draw,fill=blue!8,rounded corners=3pt,minimum width=6.7cm,minimum height=0.62cm,align=left,inner sep=3pt},
  arr/.style={-{Latex[length=2mm]},thick}
  ]
  \node[stage] (s1) {\textbf{1.}~Atualiza créditos recebidos das portas vizinhas};
  \node[stage,below=of s1] (s2) {\textbf{2.}~Lê \emph{flits} válidos dos enlaces para os \emph{buffers} FIFO};
  \node[stage,below=of s2,fill=orange!12] (s3) {\textbf{3.}~\textbf{Árbitro:} sorteia (\texttt{std::mt19937}) a ordem das 8 portas};
  \node[stage,below=of s3,fill=green!12] (s4) {\textbf{4a.}~\emph{Header}: \textbf{aloca} canal (UP adaptativo / DOWN determinístico) e \textbf{tranca}};
  \node[stage,below=of s4,fill=green!12] (s5) {\textbf{4b.}~\emph{Transmite} se há crédito; \emph{Tail} \textbf{libera} o canal};
  \node[stage,below=of s5] (s6) {\textbf{5.}~Escreve saídas de dados e devolve créditos};
  \draw[arr] (s1)--(s2); \draw[arr] (s2)--(s3); \draw[arr] (s3)--(s4);
  \draw[arr] (s4)--(s5); \draw[arr] (s5)--(s6);
\end{tikzpicture}
\caption{Pipeline lógico executado pelo processo \texttt{routing\_logic} a cada
ciclo de relógio (máquina \emph{wormhole}).}
\label{fig:fsm}
\end{figure}

\subsubsection{Árbitro Probabilístico Híbrido}
A cada ciclo, o roteador sorteia um valor em $[0,100)$ que define a ordem de
varredura das 8 portas (índices 0--3 descendentes, 4--7 ascendentes), conforme
as probabilidades da Tabela~\ref{tab:roleta}, derivadas das especificações da
SPIN. Essa aleatoriedade reordena dinamicamente as prioridades, evitando que uma
entrada perca sistematicamente a disputa por uma saída (\emph{anti-starvation}).

\begin{table}[!t]
\caption{Roleta de prioridades do árbitro}
\label{tab:roleta}
\centering
\renewcommand{\arraystretch}{1.2}
\begin{tabular}{@{}ccl@{}}
\toprule
\textbf{Prob.} & \textbf{Ordem de varredura} & \textbf{Política} \\
\midrule
50\% & 4\,5\,6\,7\,0\,1\,2\,3 & Prioridade UP (descendo da raiz) \\
30\% & 0\,1\,2\,3\,4\,5\,6\,7 & Prioridade DOWN (subindo da folha) \\
15\% & 4\,0\,5\,1\,6\,2\,7\,3 & Misto, leve vantagem UP \\
5\%  & 0\,4\,1\,5\,2\,6\,3\,7 & Misto, leve vantagem DOWN \\
\bottomrule
\end{tabular}
\end{table}

\subsubsection{Roteamento em Duas Fases}
Ao processar um \emph{header}, o roteador calcula o roteiro a partir do destino:
o roteador-folha de destino é $\lfloor \texttt{dest}/4 \rfloor$ e a porta local
de terminal é $\texttt{dest} \bmod 4$.
\begin{itemize}
  \item \textbf{Subida adaptativa:} se o roteador é de folha (\texttt{level==1})
  e o destino está em outra sub-árvore, o pacote deve subir. O roteador varre as
  4 portas ascendentes livres e escolhe a de \emph{maior número de créditos}
  disponíveis, isto é, o enlace menos congestionado.
  \item \textbf{Descida determinística:} caso contrário, a saída é unicamente
  determinada --- na raiz, pela folha de destino; na folha, pela porta do
  terminal de destino.
\end{itemize}
Uma vez escolhida, a saída é registrada em \texttt{input\_locks} e marcada como
ocupada; somente o \emph{flit-cauda} reverte essa trava, encerrando o
``túnel'' \emph{wormhole}.

\subsection{O Terminal}
O módulo \texttt{Terminal} atua como interface de rede. Quando programado por
\texttt{set\_target()}, monta uma mensagem de 4 \emph{flits}
(1~\emph{header} + 2~\emph{payload} + 1~\emph{tail}) e a injeta respeitando o
crédito disponível. Em modo de recepção, lê o enlace de chegada, imprime o
conteúdo de cada \emph{flit} e, ao receber o \emph{tail}, sinaliza a remontagem
completa do pacote, devolvendo créditos ao roteador a montante.

% ============================== RESULTADOS ==================================
\section{Resultados}
\label{sec:resultados}

\subsection{Metodologia e Cenário de Simulação}
O modelo foi compilado com \texttt{g++}~16.1 (C++17) e ligado à biblioteca
\textbf{SystemC 3.0.2 da Accellera}; os logs e temporizações apresentados nesta
seção foram \emph{extraídos de execuções concretas}, e não de estimativas
analíticas. Cada cenário de simulação é um \emph{testbench} \texttt{sc\_main}
(diretório \texttt{tests/}) que constrói a rede, aplica o relógio e o
\emph{reset}, injeta o tráfego e executa a simulação de ciclo exato --- ou seja,
\emph{os testes são a própria simulação}, acrescidos de uma verificação
automática da entrega. No cenário básico (\texttt{test\_single\_flow}):
aplica-se \emph{reset} por \SI{20}{ns}, libera-se a rede e, em $t=\SI{30}{ns}$,
configura-se o tráfego \texttt{configure\_traffic(1, 14)} --- o terminal $T_1$
envia uma mensagem \emph{wormhole} de 4 \emph{flits} ao terminal $T_{14}$; a
simulação prossegue por \SI{300}{ns}. Esse cenário é deliberadamente um
\emph{caminho longo}: origem e destino residem em sub-árvores distintas,
exigindo subida até a raiz e descida até a folha oposta. O log detalhado é
emitido sob demanda (variável de ambiente \texttt{SPIN\_TRACE}).

\subsection{Caminho Esperado}
Aplicando o endereçamento hierárquico ao destino $14$: folha de destino
$=\lfloor 14/4\rfloor = 3$ e porta de terminal $=14 \bmod 4 = 2$. A origem $T_1$
pertence à folha $\lfloor 1/4 \rfloor = 0$. Como as sub-árvores diferem
($0 \neq 3$), o pacote percorre:
\begin{equation}
T_1 \rightarrow L_0 \xrightarrow{\text{UP adapt.}} R_0
\xrightarrow{\text{DOWN det.}} L_3 \xrightarrow{\text{DOWN det.}} T_{14}.
\nonumber
\end{equation}
Em $L_0$, todas as portas ascendentes iniciam com 4 créditos; o roteamento
adaptativo seleciona a primeira de maior crédito, a porta UP~0, conectada a
$R_0$. Em $R_0$ (nível 2), a descida determinística aponta para a folha 3
(porta DOWN~3); em $L_3$, a descida final aponta para a porta DOWN~2, que serve
$T_{14}$.

\subsection{Simulação 1: Alocação do Canal pelo \emph{Header}}
O \emph{flit-cabeçalho} dispara a fase de alocação em cada salto. O excerto de
log a seguir, \emph{copiado da saída de console da execução}, evidencia a
intervenção do árbitro e o trancamento adaptativo do enlace ascendente em $L_0$:

\noindent\begin{minipage}{\columnwidth}
\begin{lstlisting}[style=log]
[ PACOTE INICIADO ] @ 30 ns | Origem: 1 Destino: 14
  [ARBITRO] Roteador Nivel 1 ID 0 | caiu em 50% (Prioridade UP) -> Atendendo Entrada 1
      -> [HW] Roteamento Adaptativo trancou porta UP 0
    >> @ 40 ns | Roteador Nivel 1 (ID 0) despachou [HEADER]  -> Porta UP 0
  [ARBITRO] Roteador Nivel 2 ID 0 | caiu em 30% (Prioridade DOWN) -> Atendendo Entrada 0
    >> @ 50 ns | Roteador Nivel 2 (ID 0) despachou [HEADER]  -> Porta DOWN 3
  [ARBITRO] Roteador Nivel 1 ID 3 | caiu em 30% (Prioridade DOWN) -> Atendendo Entrada 4
    >> @ 60 ns | Roteador Nivel 1 (ID 3) despachou [HEADER]  -> Porta DOWN 2
    <<< @ 70 ns | Terminal 14 RECEBEU [HEADER]  | Dado: >> INICIO_DA_MENSAGEM <<
\end{lstlisting}
\end{minipage}

A leitura confirma três pontos: (i)~a cada salto o árbitro registra em qual faixa
da roleta caiu (neste \emph{run}: $L_0$ em 50\% e os demais em 30\%) e qual
entrada atendeu; (ii)~em $L_0$ o roteamento adaptativo de fato trancou a porta
UP~0 (a primeira de maior crédito); e (iii)~o cabeçalho avança um salto por ciclo
de \SI{10}{ns} (\SI{40}{ns}$\to$\SI{50}{ns}$\to$\SI{60}{ns}$\to$\SI{70}{ns}),
caracterizando o avanço \emph{pipeline} típico da comutação \emph{wormhole}.

\subsection{Simulação 2: Pipeline do Corpo e Remontagem}
Após a abertura do túnel, os \emph{flits} de corpo e cauda seguem o mesmo caminho
sem nova alocação (a trava já existe), formando um \emph{pipeline sistólico}: em
um mesmo ciclo, saltos distintos transportam \emph{flits} distintos da mesma
mensagem. A Tabela~\ref{tab:pipeline} reconstrói, a partir do log de execução, a
ocupação de cada enlace ao longo do tempo.

\begin{table}[!t]
\caption{Ocupação dos enlaces por ciclo (execução real)}
\label{tab:pipeline}
\centering
\renewcommand{\arraystretch}{1.15}
\footnotesize
\begin{tabular}{@{}c|cccc@{}}
\toprule
\textbf{t (ns)} & $T_1\!\to\!L_0$ & $L_0\!\to\!R_0$ & $R_0\!\to\!L_3$ & $L_3\!\to\!T_{14}$ \\
\midrule
40  & \texttt{PAY\#1} & \texttt{HEADER} & --              & --              \\
50  & \texttt{PAY\#2} & \texttt{PAY\#1} & \texttt{HEADER} & --              \\
60  & \texttt{TAIL}   & \texttt{PAY\#2} & \texttt{PAY\#1} & \texttt{HEADER} \\
70  & --              & \texttt{TAIL}   & \texttt{PAY\#2} & \texttt{PAY\#1} \\
80  & --              & --              & \texttt{TAIL}   & \texttt{PAY\#2} \\
90  & --              & --              & --              & \texttt{TAIL}   \\
\bottomrule
\end{tabular}
\end{table}

No destino, o terminal remonta a mensagem completa, conforme o log real:

\noindent\begin{minipage}{\columnwidth}
\begin{lstlisting}[style=log]
    <<< @ 80 ns  | Terminal 14 RECEBEU [PAYLOAD] | Dado: Ola mundo :) #1
    <<< @ 90 ns  | Terminal 14 RECEBEU [PAYLOAD] | Dado: Ola mundo :) #2
    <<< @ 100 ns | Terminal 14 RECEBEU [TAIL]    | Dado: >> FIM_DA_MENSAGEM <<
        -> (Pacote completamente recebido e remontado!)
\end{lstlisting}
\end{minipage}

A ordem de chegada (\texttt{HEADER}, \texttt{PAYLOAD\#1}, \texttt{PAYLOAD\#2},
\texttt{TAIL}) é idêntica à ordem de injeção, comprovando a preservação de ordem
do \emph{wormhole} e a integridade da mensagem. A latência fim-a-fim do
cabeçalho é de 4~ciclos (\SI{40}{ns}, de $t=\SI{30}{ns}$ a $t=\SI{70}{ns}$) e o
pacote inteiro é entregue em \SI{100}{ns}.

\subsection{Verificação e Depuração}
Durante a verificação por execução, observou-se inicialmente que o terminal de
destino recebia o conteúdo do primeiro \emph{payload} duplicado, sem jamais
exibir o segundo. A investigação revelou a causa-raiz no operador de igualdade da
estrutura \texttt{Flit}: por não comparar o campo \texttt{data}, dois
\emph{payloads} de mesmo tipo, origem e destino eram considerados idênticos pelo
canal \texttt{sc\_signal}, que então suprimia a atualização do valor (semântica
padrão do SystemC, que só propaga quando \texttt{novo != atual}). A correção
consistiu em incluir o campo \texttt{data} na comparação, após o que a simulação
passou a entregar corretamente \texttt{\#1} e \texttt{\#2}, como mostra o log
acima. O episódio ilustra o valor da simulação \emph{cycle-accurate} para expor
defeitos sutis de modelagem que a inspeção estática do código não revela.

Um segundo defeito foi exposto sob tráfego concorrente: quando dois fluxos
convergiam para a mesma porta de saída, o \emph{flit-cauda} de um deles, ao
liberar o canal, permitia que o outro fluxo o readquirisse \emph{no mesmo
ciclo}, sobrescrevendo o sinal de saída e perdendo a cauda. A correção restringe
a no máximo uma escrita por porta de saída por ciclo, preservando a integridade
das mensagens sob contenção.

\subsection{Bateria de Casos de Teste}
Para assegurar que o modelo se comporta conforme especificado, foi desenvolvida
uma bateria de casos de teste autoverificáveis, integrados ao \texttt{CTest}.
Cada teste injeta um cenário, executa a simulação e \emph{compara
automaticamente} o número de pacotes remontados em cada destino com o valor
esperado, retornando falha em caso de divergência. A Tabela~\ref{tab:testes}
resume os cinco cenários; \textbf{todos passam}.

\begin{table}[!t]
\caption{Bateria de casos de teste (todos aprovados)}
\label{tab:testes}
\centering
\renewcommand{\arraystretch}{1.25}
\footnotesize
\begin{tabular}{@{}p{1.7cm}p{2.2cm}p{3.0cm}@{}}
\toprule
\textbf{Cenário} & \textbf{Fluxos} & \textbf{Propriedade verificada} \\
\midrule
Fluxo único    & $T_1\!\to\!T_{14}$        & Caminho completo (subida adaptativa + descida determinística). \\
Fluxo local    & $T_0\!\to\!T_3$           & Roteamento local na mesma folha, sem subir à raiz. \\
Convergência   & $T_0,T_8\!\to\!T_7$       & Ausência de perda de \emph{flits} sob contenção de saída. \\
Permutação     & 8 fluxos disjuntos        & $N$ fluxos concorrentes sem \emph{deadlock} nem perda. \\
Muitos-para-um & $T_0\ldots T_3\!\to\!T_{10}$ & Serialização justa de 4 fluxos para um destino. \\
\bottomrule
\end{tabular}
\end{table}

Dois cenários merecem destaque. No teste de \emph{convergência}, os dois pacotes
destinados a $T_{7}$ são entregues de forma serializada, com as caudas chegando
em \SI{100}{ns} e \SI{140}{ns} (intervalo de \SI{40}{ns}, exatamente uma
mensagem de quatro \emph{flits}), confirmando a correção do defeito de contenção.
No teste \emph{muitos-para-um}, os quatro pacotes endereçados a $T_{10}$ chegam
em \SI{100}{ns}, \SI{140}{ns}, \SI{180}{ns} e \SI{220}{ns} --- espaçamento regular que
evidencia a serialização justa e a ausência de perda mesmo sob forte disputa por
um único recurso de saída.

\subsection{Análise de Formas de Onda}
A Fig.~\ref{fig:wave} ilustra a temporização dos sinais ao longo do caminho. Em
cada salto, o sinal \texttt{valid} do enlace sobe por um ciclo quando um
\emph{flit} é despachado, e o sinal de \texttt{credit} de retorno sobe um ciclo
depois, quando o \emph{flit} é consumido a jusante, devolvendo a posição de
\emph{buffer}. O deslocamento progressivo de um ciclo entre os enlaces
$L_0\!\to\!R_0\!\to\!L_3\!\to\!T_{14}$ evidencia o avanço \emph{wormhole}.

\begin{figure}[!t]
\centering
\begin{tikzpicture}[font=\scriptsize,x=0.52cm,y=0.55cm]
  % parametros
  \def\n{9}      % numero de ciclos
  % linhas-guia verticais (ciclos)
  \foreach \c in {0,...,\n} {
    \draw[gray!25,line width=0.2pt] (\c,-0.4) -- (\c,6.4);
  }
  % CLK
  \node[anchor=east] at (-0.2,6) {clk};
  \draw[thick] (0,5.7) \foreach \c in {0,...,8} { -- (\c+0.5,5.7) -- (\c+0.5,6.3) -- (\c+1,6.3) -- (\c+1,5.7) };
  % helper: pulso em ciclo dado, na linha y
  % L0->R0 valid (ciclo 1)
  \node[anchor=east] at (-0.2,4.8) {$L_0\!\to\!R_0$.valid};
  \draw[thick,blue] (0,4.6) -- (1,4.6) -- (1,5.2) -- (2,5.2) -- (2,4.6) -- (9,4.6);
  % R0->L3 valid (ciclo 2)
  \node[anchor=east] at (-0.2,3.8) {$R_0\!\to\!L_3$.valid};
  \draw[thick,blue] (0,3.6) -- (2,3.6) -- (2,4.2) -- (3,4.2) -- (3,3.6) -- (9,3.6);
  % L3->T14 valid (ciclo 3)
  \node[anchor=east] at (-0.2,2.8) {$L_3\!\to\!T_{14}$.valid};
  \draw[thick,blue] (0,2.6) -- (3,2.6) -- (3,3.2) -- (4,3.2) -- (4,2.6) -- (9,2.6);
  % credit T14->L3 (ciclo 4)
  \node[anchor=east] at (-0.2,1.8) {$T_{14}\!\to\!L_3$.credit};
  \draw[thick,red] (0,1.6) -- (4,1.6) -- (4,2.2) -- (5,2.2) -- (5,1.6) -- (9,1.6);
  % credit L3->R0 (ciclo 5)
  \node[anchor=east] at (-0.2,0.8) {$L_3\!\to\!R_0$.credit};
  \draw[thick,red] (0,0.6) -- (5,0.6) -- (5,1.2) -- (6,1.2) -- (6,0.6) -- (9,0.6);
  % rotulos de tempo (ns reais: x=c  <->  t = 30 + 10*c)
  \foreach \c in {1,...,8} {
    \pgfmathtruncatemacro{\tns}{30+10*\c}
    \node[gray,font=\tiny] at (\c,-0.25) {\tns};
  }
  \node[gray,font=\tiny,anchor=west] at (0,-0.85) {tempo de simulação (ns), período = 10\,ns};
\end{tikzpicture}
\caption{Formas de onda do \emph{flit-header} ao longo do caminho, com os tempos
da execução real. Os pulsos \texttt{valid} (azul) propagam-se um ciclo por salto
($L_0\!\to\!R_0$ em \SI{40}{ns}, $R_0\!\to\!L_3$ em \SI{50}{ns},
$L_3\!\to\!T_{14}$ em \SI{60}{ns}); os pulsos de \texttt{credit} de retorno
(vermelho) confirmam o consumo e a liberação de \emph{buffer} a jusante.}
\label{fig:wave}
\end{figure}

\subsection{Discussão}
Os resultados validam funcionalmente os três mecanismos centrais. A
\textbf{comutação wormhole} é confirmada pela alocação no \emph{header}, pelo
fluxo em fila dos \emph{payloads} e pela liberação no \emph{tail}. O
\textbf{controle de fluxo por créditos} é confirmado pela ausência de perda de
\emph{flits} (todos os 4 chegam) e pelo retorno de créditos visível nas formas de
onda. O \textbf{roteamento adaptativo} escolheu o enlace ascendente menos
congestionado, e a \textbf{descida determinística} entregou o pacote exatamente
ao $T_{14}$. Por fim, o \textbf{árbitro probabilístico} interveio em cada salto,
registrando faixas distintas da roleta, o que ilustra a reordenação dinâmica de
prioridades projetada para mitigar \emph{starvation}.

% ============================== CONCLUSAO ===================================
\section{Conclusões}
\label{sec:conclusao}

Este trabalho apresentou um simulador \emph{cycle-accurate} em SystemC de uma
Rede-em-Chip SPIN com topologia \emph{fat-tree} quaternária de dois níveis,
contemplando 8 roteadores e 16 terminais. O modelo implementa, de forma
integrada e modular, os mecanismos que definem o desempenho de uma NoC moderna:
comutação \emph{wormhole} com alocação e liberação explícita de canais por meio
de uma tabela de travas; controle de fluxo \emph{credit-based} que assegura, por
construção, a ausência de perda de \emph{flits}; roteamento híbrido (subida
adaptativa por créditos e descida determinística por endereço hierárquico); e um
árbitro probabilístico inspirado na SPIN para mitigar \emph{starvation}.

A verificação funcional, conduzida por um cenário de tráfego de caminho longo
($T_1 \rightarrow T_{14}$), demonstrou a correta abertura do túnel
\emph{wormhole}, o avanço \emph{pipeline} de um salto por ciclo, a preservação da
ordem dos \emph{flits} e a remontagem completa da mensagem no destino, com os
créditos retornando conforme esperado. Os logs de simulação e as formas de onda
corroboram, ponto a ponto, o comportamento projetado.

% ============================== REFERENCIAS =================================
\begin{thebibliography}{9}
\bibitem{guerrier2000spin}
P.~Guerrier and A.~Greiner, ``A generic architecture for on-chip packet-switched
interconnections,'' in \emph{Proc. Design, Automation and Test in Europe (DATE)},
2000, pp.~250--256.

\bibitem{benini2002networks}
L.~Benini and G.~De~Micheli, ``Networks on chips: a new SoC paradigm,''
\emph{IEEE Computer}, vol.~35, no.~1, pp.~70--78, Jan. 2002.

\bibitem{dally2004principles}
W.~J.~Dally and B.~Towles, \emph{Principles and Practices of Interconnection
Networks}. San Francisco, CA: Morgan Kaufmann, 2004.

\bibitem{systemc}
IEEE, \emph{IEEE Standard for Standard SystemC Language Reference Manual},
IEEE Std~1666-2011, 2012.

\bibitem{adriahn2009}
T.~Bjerregaard and S.~Mahadevan, ``A survey of research and practices of
Network-on-Chip,'' \emph{ACM Computing Surveys}, vol.~38, no.~1, pp.~1--51, 2006.
\end{thebibliography}

\end{document}
